\documentclass[10pt]{article}
\usepackage[T1]{fontenc}\usepackage[utf8]{inputenc}\usepackage{textcomp}
\usepackage[scaled]{helvet}\renewcommand{\familydefault}{\sfdefault}
\DeclareUnicodeCharacter{03B5}{\ensuremath{\varepsilon}}\DeclareUnicodeCharacter{03B2}{\ensuremath{\beta}}
\DeclareUnicodeCharacter{03B7}{\ensuremath{\eta}}\DeclareUnicodeCharacter{03B1}{\ensuremath{\alpha}}
\DeclareUnicodeCharacter{03C7}{\ensuremath{\chi}}\DeclareUnicodeCharacter{2212}{\ensuremath{-}}
\DeclareUnicodeCharacter{2265}{\ensuremath{\geq}}\DeclareUnicodeCharacter{2264}{\ensuremath{\leq}}
\DeclareUnicodeCharacter{2192}{\ensuremath{\rightarrow}}
\usepackage[a4paper,margin=10mm]{geometry}
\usepackage{graphicx}
\pagestyle{empty}\setlength{\parindent}{0pt}\setlength{\parskip}{3pt}
\begin{document}
{\small Online Resource 12\hfill Fig. S1}\par\vspace{2mm}
\begin{center}\includegraphics[width=174mm,height=226mm,keepaspectratio]{/mnt/data/shahar/gwas_barley/morexV3_analysis/barley_for_publication_project/11_USED_check_low_maf_geographic_distribution_minor_allele/results/figures/Fig_carrier_origin.png}\end{center}
\vspace{-2mm}{\small \textbf{Fig. S1} Geographic origin of the accessions carrying the lead-SNP alleles and the haplotypes of the four candidate genes. \textbf{a} Genotype at the lead SNP of each of the 36 genome-wide significant loci, grouped by trait: minor-allele carrier (black), major-allele carrier (light tan) and missing call (gray). Each row is labeled with the chromosome and position (Mb) of the lead SNP, its minor-allele frequency, the direction of the minor allele's effect on the trait (+, increase; −, decrease) and, where one was carried forward, the candidate gene. \textbf{b} Haplotype groups of the four candidate genes: member of the group (black), member of another group of the same gene (light tan) and accession not assigned to a group (gray). Each row is labeled with the group, its size and its mean genotype BLUP (\% of grain; fiber / starch for the GDSL esterase/lipase). Columns are the 290 \textit{H. spontaneum} accessions, ordered by ecological region (HZ1, North–Coast; HZ2, North–Desert; HZ3, Coast–Desert transitional zones) and by sampling site\par}
\end{document}
